\bta{2023}


\begin{enumerate}
	%\renewcommand{\labelenumi}{\arabic{enumi}.}
	% A(\Alph) a(\alph) I(\Roman) i(\roman) 1(\arabic)
	%设定全局标号series=example	%引用全局变量resume=example
	%[topsep=-0.3em,parsep=-0.3em,itemsep=-0.3em,partopsep=-0.3em]
	%可使用leftmargin调整列表环境左边的空白长度 [leftmargin=0em]
	\item
	\begin{enumerate}
		%\renewcommand{\labelenumi}{\arabic{enumi}.}
		% A(\Alph) a(\alph) I(\Roman) i(\roman) 1(\arabic)
		%设定全局标号series=example	%引用全局变量resume=example
		%[topsep=-0.3em,parsep=-0.3em,itemsep=-0.3em,partopsep=-0.3em]
		%可使用leftmargin调整列表环境左边的空白长度 [leftmargin=0em]
		\item
		对于一维谐振子 $H=\frac{\hat{p}^2}{2 m}+\frac{1}{2} m \omega^2 x^2$ 有产生与湮灭算符, 分别为
		\[
		a ^{\dagger}=\frac{-i p}{\sqrt{2 m \hbar \omega}}+\sqrt{\frac{m \omega}{2 \hbar}} x  \qquad a=\frac{i p}{\sqrt{2 m \hbar \omega}}+\sqrt{\frac{m \omega}{2 \hbar}} x 
		\]
		计算 $[a, a ^{\dagger}],[a ^{\dagger}, a]$;
		\item 
		如果对于谐振子基态, $a u_0=0$, 求波函数 $u_0$;
		\item 
		把 Hamilton 量用 $a$ 与 $a ^{\dagger}$ 表示, 并求 $(a ^{\dagger})^n\left|u_0\right\rangle$ 的本征值, 并求 $u_n$ 的 本征值;
		\item 
		求 $\left\langle u_n\left|x^2\right| u_n\right\rangle$。
		
		
		
	\end{enumerate}
	
	\banswer{
	\begin{enumerate}
		\item 我们有
		\begin{align*}
			[a, a ^{\dagger}]&=[\frac{i p}{\sqrt{2 m \hbar \omega}}+\sqrt{\frac{m \omega}{2 \hbar}} x, \frac{-i p}{\sqrt{2 m \hbar \omega}}+\sqrt{\frac{m \omega}{2 \hbar}} x]\\
			&=[\frac{ip}{\sqrt{2m\hbar\omega}},\sqrt{\frac{m\omega}{2\hbar}}x]+[\sqrt{\frac{m\omega}{2\hbar}}x,\frac{-ip}{\sqrt{2m\hbar\omega}}]\\
			&=\frac{1}{2}+\frac{1}{2}\\
			&=1
		\end{align*}
		同理
		\begin{align*}
			[a^{\dagger}, a]&=-1
		\end{align*}
		
		\item 由 $a u_0=0$ 且 $p=-i\hbar\frac{d}{dx}$，得
		\begin{align*}
			a u_0&=\left(\frac{i(-i\hbar)}{\sqrt{2m\hbar\omega}}\frac{d}{dx}+\sqrt{\frac{m\omega}{2\hbar}}x\right)u_0\\
			&=\left(\frac{\hbar}{\sqrt{2m\hbar\omega}}\frac{d}{dx}+\sqrt{\frac{m\omega}{2\hbar}}x\right)u_0=0
		\end{align*}
		即
		\begin{align*}
			\frac{du_0}{dx}=-\frac{m\omega}{\hbar}x\,u_0
		\end{align*}
		解得 $u_0(x)=C\exp\!\left(-\frac{m\omega}{2\hbar}x^2\right)$。由归一化
		\begin{align*}
			\int_{-\infty}^{+\infty}|u_0|^2dx=C^2\sqrt{\frac{\pi\hbar}{m\omega}}=1
		\end{align*}
		得 $C=\left(\frac{m\omega}{\pi\hbar}\right)^{1/4}$，故
		\begin{align*}
			u_0(x)=\left(\frac{m\omega}{\pi\hbar}\right)^{1/4}\exp\!\left(-\frac{m\omega}{2\hbar}x^2\right)
		\end{align*}
		
		\item 由 $x=\sqrt{\frac{\hbar}{2m\omega}}(a+a^{\dagger})$，$p=i\sqrt{\frac{m\hbar\omega}{2}}(a^{\dagger}-a)$，代入 $H$ 得
		\begin{align*}
			H&=\frac{p^2}{2m}+\frac{1}{2}m\omega^2x^2\\
			&=\frac{\hbar\omega}{4}(a+a^{\dagger})^2-\frac{\hbar\omega}{4}(a^{\dagger}-a)^2\\
			&=\frac{\hbar\omega}{2}(aa^{\dagger}+a^{\dagger}a)\\
			&=\hbar\omega\left(a^{\dagger}a+\frac{1}{2}\right)
		\end{align*}
		令 $N=a^{\dagger}a$，则 $H=\hbar\omega(N+\frac{1}{2})$。利用 $[a,(a^{\dagger})^n]=n(a^{\dagger})^{n-1}$，有
		\begin{align*}
			N(a^{\dagger})^n|u_0\rangle&=a^{\dagger}a(a^{\dagger})^n|u_0\rangle=a^{\dagger}[(a^{\dagger})^na+n(a^{\dagger})^{n-1}]|u_0\rangle\\
			&=n(a^{\dagger})^n|u_0\rangle
		\end{align*}
		（因为 $a|u_0\rangle=0$）。故
		\begin{align*}
			H(a^{\dagger})^n|u_0\rangle=\hbar\omega\left(n+\frac{1}{2}\right)(a^{\dagger})^n|u_0\rangle
		\end{align*}
		$(a^{\dagger})^n|u_0\rangle$ 的本征值为 $\hbar\omega(n+\frac{1}{2})$。归一化后 $|u_n\rangle=\frac{1}{\sqrt{n!}}(a^{\dagger})^n|u_0\rangle$，故 $u_n$ 的本征值为
		\begin{align*}
			E_n=\hbar\omega\left(n+\frac{1}{2}\right)
		\end{align*}
		
		\item 由 $x=\sqrt{\frac{\hbar}{2m\omega}}(a+a^{\dagger})$，得
		\begin{align*}
			x^2=\frac{\hbar}{2m\omega}(a^2+aa^{\dagger}+a^{\dagger}a+a^{\dagger 2})
		\end{align*}
		由于 $a|u_n\rangle=\sqrt{n}|u_{n-1}\rangle$，$a^{\dagger}|u_n\rangle=\sqrt{n+1}|u_{n+1}\rangle$，交叉项矩阵元为零，且
		\begin{align*}
			\langle u_n|aa^{\dagger}|u_n\rangle=n+1,\quad \langle u_n|a^{\dagger}a|u_n\rangle=n
		\end{align*}
		故
		\begin{align*}
			\langle u_n|x^2|u_n\rangle=\frac{\hbar}{2m\omega}(2n+1)=\frac{\hbar}{m\omega}\left(n+\frac{1}{2}\right)
		\end{align*}
	\end{enumerate}
		
	}
	
\item 
自旋波函数可表示为 $\psi(x)=\left(\begin{array}{c}\xi(x) \\ \eta(x)\end{array}\right)$, 若 $H=p \cdot \sigma_y++m(x) \sigma_z$, 其中 $p$ 为动量, $\mathrm{m}$ 为一个函数 (已取 $\hbar$ 为$  1 $)。
\begin{enumerate}
	%\renewcommand{\labelenumi}{\arabic{enumi}.}
	% A(\Alph) a(\alph) I(\Roman) i(\roman) 1(\arabic)
	%设定全局标号series=example	%引用全局变量resume=example
	%[topsep=-0.3em,parsep=-0.3em,itemsep=-0.3em,partopsep=-0.3em]
	%可使用leftmargin调整列表环境左边的空白长度 [leftmargin=0em]
	\item
	求 Hamilton 量的坐标表象表示;
	\item 
	若 $m(x)=m$, 不随空间变化, 求动量为 $p$ 的 Hamilton 量本征态;
	\item 
	$\sigma_x$ 为泡利算符: 如果 $\sigma_x \psi(x)$ 本征值为$  1 $ , 则手征为$  1 $ , 求如果手征为$  1 $ ， $\xi(x)$ 与 $\eta(x)$ 满足什么条件;
	\item 
	如果 $m(x)=-M \tanh \left(x / x_0\right)$, 且波函数满足手征为$  1 $ , 求 Hamilton 量 本征值为$  0  $的态矢。(提示: $\int \tanh x d x=\ln \cosh x+C$)
	
	
\end{enumerate}

\banswer{
	\begin{enumerate}
		\item 坐标表示为
		\begin{align*}
			\hat{H}=-i\hbar\frac{\partial}{\partial x}\sigma_y+m(x)\sigma_z
		\end{align*}
		
		\item 当 $m(x)=m$ 时，对动量本征态 $\psi(x)=e^{ipx}\chi$，$\chi=\begin{pmatrix}\xi\\\eta\end{pmatrix}$ 满足 $(p\sigma_y+m\sigma_z)\chi=E\chi$，即
		\begin{align*}
			\begin{pmatrix} m & -ip \\ ip & -m \end{pmatrix}\begin{pmatrix}\xi\\\eta\end{pmatrix}=E\begin{pmatrix}\xi\\\eta\end{pmatrix}
		\end{align*}
		久期方程 $\det(H-EI)=E^2-(p^2+m^2)=0$，得 $E=\pm\sqrt{p^2+m^2}$。
		
		对 $E=+\sqrt{p^2+m^2}$，由第二行 $ip\xi=(E+m)\eta$，取 $\xi=E+m$，$\eta=ip$，归一化得
		\begin{align*}
			\chi_+=\frac{1}{\sqrt{2E(E+m)}}\begin{pmatrix}E+m\\ip\end{pmatrix}
		\end{align*}
		
		对 $E=-\sqrt{p^2+m^2}$，由第一行 $(m-E)\xi=ip\eta$，取 $\xi=ip$，$\eta=m-E$，归一化得
		\begin{align*}
			\chi_-=\frac{1}{\sqrt{2|E|(|E|+m)}}\begin{pmatrix}ip\\|E|+m\end{pmatrix}
		\end{align*}
		其中 $|E|=\sqrt{p^2+m^2}$。
		
		\item 手征为 $1$ 即 $\sigma_x\psi=\psi$，而
		\begin{align*}
			\sigma_x\psi=\begin{pmatrix}0&1\\1&0\end{pmatrix}\begin{pmatrix}\xi\\\eta\end{pmatrix}=\begin{pmatrix}\eta\\\xi\end{pmatrix}
		\end{align*}
		故 $\xi=\eta$。
		
		\item 手征为 $1$ 时 $\xi=\eta\equiv\phi(x)$，$E=0$ 的定态方程为
		\begin{align*}
			H\begin{pmatrix}\phi\\\phi\end{pmatrix}=\begin{pmatrix}m(x)\phi-\phi'\\\phi'-m(x)\phi\end{pmatrix}=0
		\end{align*}
		即 $\frac{d\phi}{dx}=m(x)\phi=-M\tanh\!\left(\frac{x}{x_0}\right)\phi$。分离变量并积分（利用提示）：
		\begin{align*}
			\ln\phi=-M\int\tanh\!\left(\frac{x}{x_0}\right)dx=-Mx_0\ln\cosh\!\left(\frac{x}{x_0}\right)+C
		\end{align*}
		得 $\phi(x)=A\left[\cosh\!\left(\frac{x}{x_0}\right)\right]^{-Mx_0}$。故
		\begin{align*}
			\psi(x)=A\begin{pmatrix}\left[\cosh\left(\frac{x}{x_0}\right)\right]^{-Mx_0}\\\left[\cosh\left(\frac{x}{x_0}\right)\right]^{-Mx_0}\end{pmatrix}
		\end{align*}
	\end{enumerate}
	
}



\newpage
\item 
有 HF 定理
\[
\frac{\partial E_n}{\partial \lambda}=\left\langle\frac{\partial H}{\partial \lambda}\right\rangle_n
\]
\begin{enumerate}
	%\renewcommand{\labelenumi}{\arabic{enumi}.}
	% A(\Alph) a(\alph) I(\Roman) i(\roman) 1(\arabic)
	%设定全局标号series=example	%引用全局变量resume=example
	%[topsep=-0.3em,parsep=-0.3em,itemsep=-0.3em,partopsep=-0.3em]
	%可使用leftmargin调整列表环境左边的空白长度 [leftmargin=0em]
	\item
用 HF 定理证明 WN 定理$ \langle\frac{p^2}{2 m}\rangle=\frac{1}{2}\langle x \frac{\partial V}{\partial x}\rangle_n $;
\item 
用 WN 定理求谐振子 $\left\langle\frac{p^2}{2 m}\right\rangle$;
\item 
已知氢原子 Hamilton 量为
\[
H=-\frac{\hbar^2}{2 m} \frac{d^2}{d r^2}+\frac{\hbar^2}{2 m} \frac{l(l+1)}{r^2}-\frac{e^2}{4 \pi \epsilon_0 r},  \qquad E_n= -\frac{m e^4}{32 \pi^2 \epsilon_0^2 \hbar^2 n^2}
\]
求 $\left\langle n l\left|\frac{1}{r}\right| n l\right\rangle,\left\langle n l\left|\frac{1}{r^2}\right| n l\right\rangle$.
\end{enumerate}

\banswer{
\begin{enumerate}
	\item 
我们知道哈密顿量$\hat{H}=\frac{\hat{p}^2}{2m}+V(\hat{x},t)=-\frac{\hbar^2}{2m}\frac{d^2}{dx^2}+V(\hat{x},t)$,可以令$x=\lambda x$，所以得到$$\hat{H}=-\frac{1}{\lambda^2}\frac{\hbar^2}{2m}\frac{d^2}{dx^2}+V(\lambda\hat{x},t)$$
根据HF定理
\begin{align*}
\langle\frac{\partial H}{\partial\lambda}\rangle_n&=\frac{\partial E_n}{\partial\lambda}=0\\
&=\langle-\frac{2}{\lambda^3}\frac{\hat{p}^2}{2m}+x\frac{dV(\lambda\hat{x},t)}{d(\lambda x)}\rangle_n
\end{align*}
所以我们可以得到
\begin{align*}
	\langle\frac{1}{\lambda^3}\frac{\hat{p}^2}{2m}\rangle_n=\frac{1}{2}\langle x\frac{dV(\lambda\hat{x},t)}{d(\lambda x)}\rangle_n
\end{align*}
将$\lambda=1$带入可得
\begin{align*}
	\langle\frac{p^2}{2 m}\rangle_n=\frac{1}{2}\langle x \frac{\partial V}{\partial x}\rangle_n
\end{align*}
	\item 
对于谐振子，可以知道它的势能为$V(\hat{x})=\frac{1}{2}m\omega^2\hat{x}^2$，所以根据(1)中的结论可知
\begin{align*}
	\langle\frac{p^2}{2 m}\rangle_n&=\frac{1}{2}\langle x \frac{\partial V}{\partial x}\rangle_n\\
									&=\frac{1}{2}\langle m\omega\hat{x}^2\rangle_n 
\end{align*}
而
\begin{align*}
	E=\langle\frac{p^2}{2 m}\rangle_n+\frac{1}{2}\langle m\omega\hat{x}^2\rangle_n=\hbar\omega(n+\frac{1}{2})
\end{align*}
所以我们可以得到
\begin{align*}
	E=\frac{1}{2}\hbar\omega(n+\frac{1}{2})
\end{align*}
	\item 
根据HF定理，我们选择参数$\lambda=\epsilon_0$(这里选择$e$或者$\pi$都可以)，可以得到
\begin{align*}
	\frac{\partial E_n}{\partial \epsilon_0}=\frac{me^4}{16\pi^2\epsilon_0^3\hbar^2n^2}=\frac{e^2}{4\pi\epsilon_0^2}\langle nl|\frac{1}{r}|nl\rangle
\end{align*}
\begin{align*}
	\langle nl|\frac{1}{r}|nl\rangle=\frac{me^2}{4\pi\epsilon_0\hbar^2n^2}
\end{align*}
对于$\langle\frac{1}{r^2} \rangle$,由于$n=j_0+l+1$,所以我们选取参数$\lambda=l$可以得到
\begin{align*}
	\frac{\partial E_n}{\partial l}=\frac{\partial E_n}{\partial n}=\frac{me^4}{16\pi^2\epsilon_0^2\hbar^2n^3}=\frac{\hbar^2(2l+1)}{2m}\langle\frac{1}{r^2} \rangle
\end{align*}
\begin{align*}
	\langle\frac{1}{r^2}\rangle=\frac{m^2e^4}{8\pi^2\epsilon_0^2\hbar^4n^3(2l+1)}
\end{align*}
\end{enumerate}
}


\item 
$\mathrm{J}$ 为单电子总角动量, $l$ 为轨道角动量, $s$ 为自旋角动量, $\vec{\sigma} \cdot l$ 为自旋--轨道 耦合项。
\begin{enumerate}
	%\renewcommand{\labelenumi}{\arabic{enumi}.}
	% A(\Alph) a(\alph) I(\Roman) i(\roman) 1(\arabic)
	%设定全局标号series=example	%引用全局变量resume=example
	%[topsep=-0.3em,parsep=-0.3em,itemsep=-0.3em,partopsep=-0.3em]
	%可使用leftmargin调整列表环境左边的空白长度 [leftmargin=0em]
	\item
	求一组对易力学量完全集;
\item 
求 $\mathrm{J}$ 的本征值, 以及 $\vec{\sigma} \cdot l$ 的本征值;
\item 
求 $\vec{\sigma} \cdot l$ 的本征态。
	
	
	
\end{enumerate}

\banswer{
\begin{enumerate}
	\item 由 $[J_i, l^2]=0$，$[J_i, s^2]=0$，$[J_i, J^2]=0$，$[J^2, J_z]=0$ 可知一组对易力学量完全集为
	\begin{align*}
		\{l^2, s^2, J^2, J_z\}
	\end{align*}
	其中 $s^2=\frac{3}{4}\hbar^2$ 为常数，也可取 $\{H, J^2, J_z, l^2\}$。
	
	\item 总角动量 $J=l+s=l+\frac{1}{2}\sigma$，由角动量耦合规则
	\begin{align*}
		j=l+\frac{1}{2}\quad\text{或}\quad j=l-\frac{1}{2}\quad(l\neq 0)
	\end{align*}
	$J^2$ 的本征值为 $j(j+1)\hbar^2$。
	
	利用 $(\vec{\sigma}\cdot\vec{l})^2=\vec{l}^2-\vec{\sigma}\cdot\vec{l}$（因 $(\vec{A}\cdot\vec{\sigma})(\vec{B}\cdot\vec{\sigma})=\vec{A}\cdot\vec{B}+i(\vec{A}\times\vec{B})\cdot\vec{\sigma}$，令 $\vec{A}=\vec{B}=\vec{l}$），设 $\vec{\sigma}\cdot\vec{l}$ 的本征值为 $\lambda$，则
	\begin{align*}
		\lambda^2+\lambda-l(l+1)=0
	\end{align*}
	解得 $\lambda=l$ 或 $\lambda=-(l+1)$。
	
	对应关系：当 $j=l+\frac{1}{2}$ 时，$\vec{\sigma}\cdot\vec{l}=l$；当 $j=l-\frac{1}{2}$ 时，$\vec{\sigma}\cdot\vec{l}=-(l+1)$。
	
	\item 本征态即耦合表象基矢 $|j, m_j, l, s\rangle$。用 $s=\frac{1}{2}$ 的 CG 系数：
	
	对 $j=l+\frac{1}{2}$：
	\begin{align*}
		|l+\tfrac{1}{2}, m_j\rangle&=\sqrt{\frac{l+m_j+\frac{1}{2}}{2l+1}}|l, m_j-\tfrac{1}{2}\rangle|+\rangle\\
		&\quad+\sqrt{\frac{l-m_j+\frac{1}{2}}{2l+1}}|l, m_j+\tfrac{1}{2}\rangle|-\rangle
	\end{align*}
	
	对 $j=l-\frac{1}{2}$：
	\begin{align*}
		|l-\tfrac{1}{2}, m_j\rangle&=-\sqrt{\frac{l-m_j+\frac{1}{2}}{2l+1}}|l, m_j-\tfrac{1}{2}\rangle|+\rangle\\
		&\quad+\sqrt{\frac{l+m_j+\frac{1}{2}}{2l+1}}|l, m_j+\tfrac{1}{2}\rangle|-\rangle
	\end{align*}
	其中 $|+\rangle=\begin{pmatrix}1\\0\end{pmatrix}$，$|-\rangle=\begin{pmatrix}0\\1\end{pmatrix}$。此即 $\vec{\sigma}\cdot\vec{l}$ 的本征态。
\end{enumerate}
}

\item 
考虑无限深势阱
\[
V(x)=\left\{\begin{aligned}
	0,&& -\dfrac{L}{2}<x<\dfrac{L}{2} \\ 
	+\infty, && \quad x<-\dfrac{L}{2}, x>\dfrac{L}{2}\end{aligned}\right.
\]
所加微扰为 $a \sin (2 \pi x / L)$, 微扰所在区域是 $-L / 2<x<L / 2$。
\begin{enumerate}
	%\renewcommand{\labelenumi}{\arabic{enumi}.}
	% A(\Alph) a(\alph) I(\Roman) i(\roman) 1(\arabic)
	%设定全局标号series=example	%引用全局变量resume=example
	%[topsep=-0.3em,parsep=-0.3em,itemsep=-0.3em,partopsep=-0.3em]
	%可使用leftmargin调整列表环境左边的空白长度 [leftmargin=0em]
	\item
	求 $a=0$ 时（无微扰时）本征态;
\item 
求 $E_n^1$ (有微扰时);
\item 
求 $\left\langle\psi_{2 n}^0 \mid \psi_{2 m}^1\right\rangle,\left\langle\psi_{2 n}^0 \mid \psi_{2 m+1}^1\right\rangle$ 和 $\left\langle\psi_{2 n+1}^0 \mid \psi_{2 m+1}^1\right\rangle$ 。
	
	
	
\end{enumerate}
提示：积化和差公式
\[
\begin{aligned}
	& \sin a \cos b=\frac{1}{2}[\sin (a+b)+\sin (a-b)] 
	& \cos a \sin b=\frac{1}{2}[\sin (a+b)-\sin (a-b)] \\
	& \cos a \cos b=\frac{1}{2}[\cos (a+b)+\cos (a-b)] 
	& \sin a \sin b=-\frac{1}{2}[\cos (a+b)-\cos (a-b)]
\end{aligned}
\]


\banswer{
\begin{enumerate}
	\item 无微扰时，势阱范围 $[-L/2, L/2]$，能量本征值为
	\begin{align*}
		E_n^0=\frac{n^2\pi^2\hbar^2}{2mL^2},\quad n=1,2,3,\ldots
	\end{align*}
	本征函数按宇称分为：
	\begin{align*}
		\psi_n^0(x)=\sqrt{\frac{2}{L}}\begin{cases}
			\cos\!\left(\dfrac{n\pi x}{L}\right), & n=1,3,5,\ldots\quad(\text{偶宇称})\\[6pt]
			\sin\!\left(\dfrac{n\pi x}{L}\right), & n=2,4,6,\ldots\quad(\text{奇宇称})
		\end{cases}
	\end{align*}
	
	\item 一级能量修正为
	\begin{align*}
		E_n^{(1)}=\langle\psi_n^0|H'|\psi_n^0\rangle=a\int_{-L/2}^{L/2}|\psi_n^0(x)|^2\sin\!\left(\frac{2\pi x}{L}\right)dx
	\end{align*}
	由于 $|\psi_n^0|^2$ 是偶函数，$\sin(2\pi x/L)$ 是奇函数，被积函数为奇函数，在对称区间积分为零，故
	\begin{align*}
		E_n^{(1)}=0
	\end{align*}
	
	\item 非简并微扰论的波函数一级修正为
	\begin{align*}
		|\psi_n^1\rangle=\sum_{k\neq n}\frac{\langle\psi_k^0|H'|\psi_n^0\rangle}{E_n^0-E_k^0}|\psi_k^0\rangle
	\end{align*}
	因此
	\begin{align*}
		\langle\psi_m^0|\psi_n^1\rangle=\frac{\langle\psi_m^0|H'|\psi_n^0\rangle}{E_n^0-E_m^0}\quad(m\neq n)
	\end{align*}
	
	对 $\langle\psi_{2n}^0|\psi_{2m}^1\rangle$：$\psi_{2n}^0$ 与 $\psi_{2m}^0$ 均为奇函数，$H'$ 为奇函数，被积函数奇$\times$奇$\times$奇$=$奇，积分为零，故
	\begin{align*}
		\langle\psi_{2n}^0|\psi_{2m}^1\rangle=0
	\end{align*}
	
	对 $\langle\psi_{2n}^0|\psi_{2m+1}^1\rangle$：$\psi_{2n}^0$ 奇，$\psi_{2m+1}^0$ 偶，$H'$ 奇，被积函数为偶函数，一般非零：
	\begin{align*}
		&\langle\psi_{2n}^0|H'|\psi_{2m+1}^0\rangle\\
		&\quad=a\frac{2}{L}\int_{-L/2}^{L/2}\sin\!\left(\frac{2n\pi x}{L}\right)\sin\!\left(\frac{2\pi x}{L}\right)\cos\!\left(\frac{(2m+1)\pi x}{L}\right)dx
	\end{align*}
	利用积化和差 $\sin A\sin B=-\frac{1}{2}[\cos(A+B)-\cos(A-B)]$，得
	\begin{align*}
		&\sin\!\left(\frac{2n\pi x}{L}\right)\sin\!\left(\frac{2\pi x}{L}\right)\\
		&\quad=-\frac{1}{2}\left[\cos\!\left(\frac{2(n+1)\pi x}{L}\right)-\cos\!\left(\frac{2(n-1)\pi x}{L}\right)\right]
	\end{align*}
	再与 $\cos\!\left(\frac{(2m+1)\pi x}{L}\right)$ 相乘后用 $\cos A\cos B=\frac{1}{2}[\cos(A+B)+\cos(A-B)]$ 展开，各余弦积分在 $[-L/2,L/2]$ 上仅当频率为零时才非零，可逐次算出。
	
	故
	\begin{align*}
		\langle\psi_{2n}^0|\psi_{2m+1}^1\rangle=\frac{\langle\psi_{2n}^0|H'|\psi_{2m+1}^0\rangle}{E_{2m+1}^0-E_{2n}^0}
	\end{align*}
	
	对 $\langle\psi_{2n+1}^0|\psi_{2m+1}^1\rangle$：$\psi_{2n+1}^0$ 与 $\psi_{2m+1}^0$ 均为偶函数，$H'$ 为奇函数，被积函数偶$\times$奇$\times$偶$=$奇，积分为零，故
	\begin{align*}
		\langle\psi_{2n+1}^0|\psi_{2m+1}^1\rangle=0
	\end{align*}
\end{enumerate}
}



\end{enumerate}

